\section{Results and Discussion}

\subsection{Model Performance}

The final production model (Logistic Regression) achieves the following metrics on the
held-out 20\% test split (1,409 samples):

\begin{itemize}[leftmargin=*]
    \item \textbf{Accuracy}: 80.7\%---correctly classifies the majority of customers.
    \item \textbf{Precision}: 65.8\%---of customers predicted to churn, about two-thirds
        actually do.
    \item \textbf{Recall}: 56.6\%---identifies more than half of actual churners.
    \item \textbf{F1-Score}: 0.609---balances precision and recall effectively.
    \item \textbf{ROC-AUC}: 0.842---strong discriminative ability between churners and
        non-churners.
\end{itemize}

The confusion matrix for the test set shows 925 true negatives, 212 true positives,
110 false positives, and 162 false negatives. The model is more conservative in its
churn predictions, which is appropriate for a retention system where false positives
(triggering retention outreach for a customer who would have stayed) are less costly
than false negatives (missing a customer who leaves).

\subsection{Qualitative Agent Evaluation}

To evaluate the agent's effectiveness beyond raw metrics, we tested it on two
representative customer profiles.

\subsubsection{High-Risk Customer Profile}

\begin{lstlisting}[caption=High-Risk customer input features.]
senior: No, tenure: 2, monthly: $90, total: $180
gender: Female, partner: No, dependents: No
phone: Yes, multilines: No, internet: Fiber optic
online_sec: No, online_bkp: No, device_prot: No, tech_sup: No
streaming_tv: Yes, streaming_movies: Yes
contract: Month-to-month, paperless: Yes
payment: Electronic check
\end{lstlisting}

The ML model predicts a churn probability of approximately 85--90\% with a ``High'' risk
classification. The agent's risk factor identification correctly flags: month-to-month
contract, short tenure (2 months), fiber optic subscriber, electronic check payment,
high monthly charges (\$90), and no value-add services.

The LLM-generated explanation correctly references the customer's short tenure, high
monthly charges, lack of online security and tech support, and month-to-month contract
as key risk indicators. The tone is professional and avoids ML jargon.

The three recommendations are personalized, referencing specific features of the
customer's profile and suggesting actionable steps such as contract upgrade offers,
auto-pay enrollment incentives, and free security add-on trials.

\subsubsection{Loyal Customer Profile}

\begin{lstlisting}[caption=Loyal customer input features.]
senior: No, tenure: 60, monthly: $55, total: $3300
gender: Male, partner: Yes, dependents: Yes
phone: Yes, multilines: Yes, internet: DSL
online_sec: Yes, online_bkp: Yes, device_prot: Yes, tech_sup: Yes
streaming_tv: No, streaming_movies: No
contract: Two year, paperless: No
payment: Bank transfer (automatic)
\end{lstlisting}

The model predicts a churn probability under 10\% with a ``Low'' risk classification.
The agent correctly identifies few or no significant risk factors. The LLM explanation
acknowledges the customer's strong retention indicators (long tenure, two-year contract,
comprehensive service bundle) and frames the output as a retention maintenance
recommendation rather than a crisis intervention.

\subsection{RAG Retrieval Quality}

For the High-Risk profile, the retriever returns chunks covering month-to-month contract
strategies, electronic check payment retention, and short-tenure onboarding strategies---all
directly relevant to the customer's identified risk factors. The retrieved context
meaningfully informs the LLM's recommendations.

For the Loyal profile, the retriever returns general retention principles, which the LLM
appropriately uses to suggest maintenance actions rather than aggressive retention offers.

\subsection{Limitations}

\begin{enumerate}[leftmargin=*]
    \item \textbf{Hallucination risk}: The LLM may occasionally generate recommendations
        not grounded in the retrieved context. Prompt engineering mitigates but does not
        eliminate this risk.
    \item \textbf{Static knowledge base}: The retention strategies document is manually
        curated and does not update automatically based on real-world outcome data.
    \item \textbf{Single-telecom training data}: The model is trained exclusively on the
        IBM Telco dataset; generalization to other telecom datasets or industries is untested.
    \item \textbf{LLM latency}: The full agent pipeline takes 5--15 seconds, primarily due
        to three sequential LLM calls. This is acceptable for a decision-support tool but
        not for real-time use.
    \item \textbf{No feedback loop}: The system does not track whether its recommendations
        actually prevent churn, limiting its ability to improve over time.
\end{enumerate}
